\chapter{Yêu cầu phi chức năng}

\section{Giới thiệu}

Bên cạnh các yêu cầu chức năng, hệ thống AITasker cần đáp ứng các yêu cầu phi
chức năng nhằm đảm bảo chất lượng phần mềm, hiệu năng, độ tin cậy, tính bảo
mật và khả năng mở rộng. Các yêu cầu này đóng vai trò quan trọng trong việc
đánh giá mức độ hoàn thiện của hệ thống sau khi triển khai.

Các yêu cầu phi chức năng được xây dựng dựa trên đặc điểm của hệ thống kết nối
doanh nghiệp và chuyên gia AI, đồng thời phù hợp với kiến trúc sử dụng ReactJS,
FastAPI, PostgreSQL và các dịch vụ AI tích hợp.

%====================================================

\section{Yêu cầu về hiệu năng}

Hệ thống phải đáp ứng các yêu cầu hiệu năng sau:

\begin{itemize}

\item Thời gian phản hồi trung bình của các API không vượt quá 2 giây.

\item Chức năng đăng nhập hoàn thành trong thời gian không quá 3 giây.

\item Chức năng tạo dự án hoàn thành trong thời gian không quá 5 giây.

\item Chức năng gửi đề xuất hoàn thành trong thời gian không quá 3 giây.

\item Chức năng thanh toán Milestone hoàn thành trong thời gian không quá 5 giây.

\item Dashboard hiển thị dữ liệu trong thời gian không quá 5 giây.

\item Hệ thống hỗ trợ tối thiểu 500 người dùng hoạt động đồng thời.

\item Hệ thống có khả năng xử lý tối thiểu 100 yêu cầu mỗi giây.

\item Truy vấn cơ sở dữ liệu thông thường không vượt quá 1 giây.

\end{itemize}

%====================================================

\section{Yêu cầu về bảo mật}

Do hệ thống quản lý tài khoản, hợp đồng và thanh toán nên các yêu cầu bảo mật
được đặt ở mức cao.

\begin{itemize}

\item Người dùng phải được xác thực bằng JWT.

\item Mật khẩu được mã hóa bằng BCrypt trước khi lưu vào cơ sở dữ liệu.

\item Không lưu mật khẩu dưới dạng văn bản thuần.

\item Token phải có thời gian hết hạn.

\item Chỉ người dùng đã đăng nhập mới được truy cập các API nghiệp vụ.

\item Mỗi API phải kiểm tra quyền truy cập theo vai trò.

\item Chỉ Enterprise được phép tạo Project.

\item Chỉ Expert được phép gửi Proposal.

\item Chỉ Administrator được phép quản trị hệ thống.

\item Dữ liệu đầu vào phải được kiểm tra để hạn chế SQL Injection.

\item Dữ liệu đầu vào phải được kiểm tra nhằm hạn chế Cross Site Scripting (XSS).

\item Hệ thống hỗ trợ truyền dữ liệu thông qua HTTPS khi triển khai thực tế.

\end{itemize}

%====================================================

\section{Yêu cầu về độ tin cậy}

Hệ thống cần đảm bảo:

\begin{itemize}

\item Không làm mất dữ liệu khi xảy ra lỗi.

\item Mọi thao tác ghi dữ liệu được thực hiện bằng Transaction.

\item Dữ liệu luôn đảm bảo tính toàn vẹn.

\item Có khả năng phục hồi sau khi xảy ra sự cố.

\item Không tạo dữ liệu trùng lặp.

\item Ghi nhận nhật ký lỗi phục vụ bảo trì.

\end{itemize}

%====================================================

\section{Yêu cầu về tính sẵn sàng}

Hệ thống phải đáp ứng:

\begin{itemize}

\item Hoạt động liên tục 24 giờ mỗi ngày.

\item Thời gian sẵn sàng tối thiểu đạt 99\%.

\item Có khả năng khởi động lại sau khi máy chủ gặp sự cố.

\item Không làm gián đoạn các giao dịch đang xử lý.

\item Người dùng có thể truy cập hệ thống mọi lúc thông qua trình duyệt Web.

\end{itemize}

%====================================================

\section{Yêu cầu về khả năng mở rộng}

Hệ thống phải hỗ trợ:

\begin{itemize}

\item Gia tăng số lượng người dùng.

\item Gia tăng số lượng doanh nghiệp.

\item Gia tăng số lượng chuyên gia AI.

\item Gia tăng số lượng dự án.

\item Gia tăng số lượng hợp đồng.

\item Mở rộng cơ sở dữ liệu mà không ảnh hưởng tới hệ thống.

\item Dễ dàng tích hợp thêm các dịch vụ AI mới.

\end{itemize}

%====================================================

\section{Yêu cầu về khả năng sử dụng}

Hệ thống phải đảm bảo:

\begin{itemize}

\item Giao diện trực quan.

\item Dễ sử dụng đối với người mới.

\item Thiết kế Responsive.

\item Hỗ trợ nhiều trình duyệt.

\item Thông báo lỗi rõ ràng.

\item Điều hướng đơn giản.

\item Tìm kiếm nhanh các dự án và chuyên gia.

\item Thời gian học sử dụng hệ thống ngắn.

\end{itemize}

%====================================================

\section{Yêu cầu về khả năng bảo trì}

Để thuận tiện trong việc nâng cấp hệ thống, phần mềm phải đáp ứng:

\begin{itemize}

\item Áp dụng kiến trúc nhiều tầng.

\item Tuân thủ RESTful API.

\item Tách biệt Frontend và Backend.

\item Sử dụng SQLAlchemy ORM.

\item Có cấu trúc thư mục rõ ràng.

\item Có tài liệu API.

\item Có Logging.

\item Có Exception Handling.

\item Dễ dàng bổ sung module mới.

\end{itemize}

%====================================================

\section{Yêu cầu về khả năng tương thích}

Hệ thống hỗ trợ:

\begin{itemize}

\item Google Chrome.

\item Microsoft Edge.

\item Mozilla Firefox.

\item Windows.

\item Linux.

\item macOS.

\item PostgreSQL.

\item Python 3.12.

\item NodeJS.

\end{itemize}

%====================================================

\section{Yêu cầu về sao lưu và phục hồi}

Hệ thống cần hỗ trợ:

\begin{itemize}

\item Sao lưu dữ liệu định kỳ.

\item Khôi phục dữ liệu khi xảy ra sự cố.

\item Lưu nhật ký hệ thống.

\item Theo dõi lịch sử đăng nhập.

\item Theo dõi lịch sử thay đổi dữ liệu.

\item Khôi phục cơ sở dữ liệu từ bản sao lưu.

\end{itemize}

%====================================================

\section{Yêu cầu về tích hợp}

Hệ thống hỗ trợ tích hợp với:

\begin{itemize}

\item Google Gemini API.

\item Email Service.

\item JWT Authentication.

\item PostgreSQL.

\item Azure Cloud Storage.

\item RESTful API.

\end{itemize}

%====================================================

\section{Ràng buộc kỹ thuật}

Bảng~\ref{tab:technology} trình bày các công nghệ chính được sử dụng trong hệ
thống.

\begin{table}[H]
\centering
\caption{Các công nghệ sử dụng trong hệ thống}
\label{tab:technology}

\begin{tabular}{|p{6cm}|p{8cm}|}
\hline
\textbf{Thành phần} & \textbf{Công nghệ} \\ \hline

Frontend & ReactJS \\ \hline

Backend & FastAPI \\ \hline

Ngôn ngữ lập trình & Python 3.12 \\ \hline

Cơ sở dữ liệu & PostgreSQL \\ \hline

ORM & SQLAlchemy \\ \hline

Authentication & JWT \\ \hline

Password Hashing & BCrypt \\ \hline

AI Platform & Google Gemini API \\ \hline

Cloud Deployment & Microsoft Azure \\ \hline

API Architecture & RESTful API \\ \hline

\end{tabular}

\end{table}

%====================================================

\section{Tiêu chí chấp nhận hệ thống}

Hệ thống được xem là đáp ứng yêu cầu khi thỏa mãn các điều kiện sau:

\begin{enumerate}

\item Người dùng có thể đăng ký và đăng nhập thành công.

\item Enterprise có thể tạo Project.

\item Expert có thể gửi Proposal.

\item Enterprise có thể chấp nhận Proposal.

\item Contract được tạo đúng quy trình.

\item Milestone được quản lý đúng trạng thái.

\item Thanh toán được thực hiện thành công.

\item Người dùng có thể đánh giá sau khi hợp đồng hoàn thành.

\item Dashboard hiển thị đúng dữ liệu.

\item JWT xác thực chính xác.

\item Phân quyền theo Role hoạt động đúng.

\item Hệ thống vận hành ổn định trên các trình duyệt được hỗ trợ.

\end{enumerate}

%====================================================

\section{Kết luận}

Các yêu cầu phi chức năng là cơ sở để đánh giá chất lượng của hệ thống
AITasker sau khi triển khai. Việc đáp ứng đầy đủ các yêu cầu này sẽ giúp hệ
thống hoạt động ổn định, bảo mật, có khả năng mở rộng và đáp ứng nhu cầu sử
dụng của doanh nghiệp cũng như chuyên gia AI trong thực tế.